\begin{explanationbox}
\textbf{\textcolor{CExplanation}{一句话直觉}}

凹函数图像上任意两点连弦位于曲线下方，因此函数值的算术平均不超过平均自变量的函数值。

\medskip
\textbf{\textcolor{CExplanation}{核心拆解}}
\begin{description}[style=nextline, leftmargin=2.8em, labelsep=0.5em, itemsep=0.45em]
    \item[\textbf{要点一（凹函数定义）：}] 对于凹函数 $f$，有 $f\left(\frac{x+y}{2}\right) \geq \frac{f(x)+f(y)}{2}$。
    \item[\textbf{要点二（Jensen不等式）：}] 对于凹函数，函数值的加权平均不超过自变量加权平均的函数值。
\end{description}

\medskip
\textbf{\textcolor{CExplanation}{几何本质（弦在曲线下）}}

在区间 $(0,\frac{\pi}{2})$ 上，余弦函数图像是凹的，连接图像上两点的弦位于曲线下方。

\medskip
\textbf{\textcolor{CExplanation}{代数意义（平均不等式）}}

不等式反映了算术平均与函数值的关系：对于凹函数，函数值的算术平均不超过自变量的算术平均的函数值。

\medskip
\textbf{\textcolor{CExplanation}{考点价值（不等式证明）}}
\begin{itemize}[leftmargin=*, itemsep=0.5em, topsep=0.4em]
    \item \textbf{考法一（直接应用）：} 证明两个角的余弦不等式。
    \item \textbf{考法二（三角形极值）：} 求锐角三角形余弦和的最大值。
\end{itemize}

\medskip
\textbf{\textcolor{CExplanation}{顿悟点}}

凹性导致“平均函数值”小于“函数平均值”，但注意方向：$\frac{f(x)+f(y)}{2} \leq f\left(\frac{x+y}{2}\right)$ 对应凹函数。

\medskip
\textbf{\textcolor{CExplanation}{使用场景}}
\begin{itemize}[leftmargin=*, itemsep=0.5em, topsep=0.4em]
    \item \textbf{场景一（证明不等式）：} 当涉及余弦函数的不等式时，可考虑凹性放缩。
    \item \textbf{场景二（优化问题）：} 在锐角三角形中求三角函数和的极值。
\end{itemize}
\end{explanationbox}
